\documentclass{article}
\usepackage{amsmath,amssymb,amsthm}
\usepackage{algorithm}
\usepackage[noend]{algpseudocode}
\usepackage{hyperref}
\newtheorem{theorem}{Theorem}
\newtheorem{lemma}{Lemma}
\newcommand{\prox}{\operatorname{prox}}
\newcommand{\R}{\mathbb{R}}

\begin{document}

\section*{Algorithm Name}
\textbf{Proximal Accelerated Gradient Method (PAGM)}  
A first‑order method for minimizing a composite convex objective  
\[
\min_{x\in\R^{d}} \; F(x)\;:=\; f(x)+g(x),
\]
where $f$ is smooth (Lipschitz gradient) and $g$ is convex, possibly nonsmooth.  
The algorithm applies Nesterov’s momentum to the \emph{gradient input} before the proximal step.

\section*{Mathematical Formulation}
Let a stepsize $\eta>0$ be given (typically $\eta\le 1/L$).  
Define sequences $\{x_k\}_{k\ge 0}$, $\{y_k\}_{k\ge 0}$ and $\{t_k\}_{k\ge 0}$ by  

\[
\begin{aligned}
x_0 &\in\R^{d},\qquad x_{-1}=x_0,\\[4pt]
t_0 &=1, \\[4pt]
\text{for }k\ge 0:\qquad 
y_k &= x_k+\beta_k\,(x_k-x_{k-1}),\qquad 
\beta_k:=\frac{t_{k-1}-1}{t_k},\\[4pt]
x_{k+1} &= \prox_{\eta g}\!\bigl(y_k-\eta\nabla f(y_k)\bigr),\\[4pt]
t_{k+1} &= \frac{1+\sqrt{1+4t_k^{2}}}{2}.
\end{aligned}
\]
The proximal operator is
\[
\prox_{\eta g}(z):=\arg\min_{u\in\R^{d}}\Bigl\{g(u)+\frac{1}{2\eta}\|u-z\|^{2}\Bigr\}.
\]

\section*{Pseudocode}
\begin{algorithm}[H]
\caption{Proximal Accelerated Gradient Method}
\begin{algorithmic}[1]
\Require stepsize $\eta>0$, initial point $x_0\in\R^{d}$, max iterations $K$
\State $x_{-1}\gets x_0$, $t_0\gets 1$
\For{$k=0,\dots,K-1$}
    \State $\beta_k\gets (t_{k-1}-1)/t_k$
    \State $y_k\gets x_k+\beta_k\,(x_k-x_{k-1})$
    \State $g_k\gets \nabla f(y_k)$
    \State $x_{k+1}\gets \prox_{\eta g}\!\bigl(y_k-\eta g_k\bigr)$
    \State $t_{k+1}\gets \frac{1+\sqrt{1+4t_k^{2}}}{2}$
\EndFor
\State \textbf{return} $x_{K}$
\end{algorithmic}
\end{algorithm}

\section*{Dry Run Example}
Minimize $F(w)=f(w)+g(w)$ with  
\[
f(w)=(w-3)^{2},\qquad g(w)=0\;(\text{so }\prox_{\eta g}= \operatorname{id}),
\]
stepsize $\eta=0.1$, start $w_0=0$, run three iterations.

\[
\begin{array}{c|c|c|c|c}
k & t_k & \beta_k & y_k & w_{k+1}=y_k-\eta\nabla f(y_k)\\ \hline
0 & 1 & 0 & y_0=w_0=0 & \nabla f(0)= -6\;\Rightarrow\; w_1=0-0.1(-6)=0.6\\
1 & \frac{1+\sqrt{5}}{2}\approx1.618 & \frac{t_0-1}{t_1}=0 & y_1=w_1=0.6 & \nabla f(0.6)=2(0.6-3)=-4.8\\
  &  &  &  & w_2=0.6-0.1(-4.8)=1.08\\
2 & \frac{1+\sqrt{1+4t_1^{2}}}{2}\approx2.193 & \beta_2=\frac{t_1-1}{t_2}\approx0.281\\
  &  &  & y_2 = w_2+\beta_2(w_2-w_1)=1.08+0.281(1.08-0.6)\approx1.215\\
  &  &  & \nabla f(y_2)=2(1.215-3)=-3.57\\
  &  &  & w_3= y_2-0.1(-3.57)\approx1.572
\end{array}
\]

Objective values: $F(w_0)=9$, $F(w_1)=5.76$, $F(w_2)=3.6864$, $F(w_3)=2.838\,$ (monotonically decreasing).

\newpage
\section*{Assumptions}
\begin{enumerate}
\item \textbf{Smooth part.} $f:\R^{d}\to\R$ is convex and differentiable with $L$‑Lipschitz continuous gradient:
\[
\|\nabla f(x)-\nabla f(y)\|\le L\|x-y\|,\qquad\forall x,y\in\R^{d}. \tag{A1}
\]
\item \textbf{Nonsmooth part.} $g:\R^{d}\to\R\cup\{+\infty\}$ is proper, closed and convex. Its proximal operator is well defined for any $\eta>0$. \tag{A2}
\item \textbf{Stepsize.} $\eta\in(0,1/L]$. \tag{A3}
\item \textbf{Initialization.} $x_{-1}=x_0\in\operatorname{dom}(F)$. \tag{A4}
\end{enumerate}

\section*{Main Convergence Theorem}
\begin{theorem}[Accelerated $O(1/k^{2})$ rate]\label{thm:rate}
Under assumptions (A1)–(A4) the sequence $\{x_k\}$ generated by PAGM satisfies for all $k\ge 1$
\[
F(x_k)-F(x^\star)\;\le\;\frac{2\|x_0-x^\star\|^{2}}{\eta\,k^{2}},
\]
where $x^\star\in\arg\min F$.
\end{theorem}

\section*{Theoretical Properties}
\begin{itemize}
\item \textbf{Stability.} The recurrence for $t_k$ guarantees $\beta_k\in[0,1)$, thus the extrapolation never overshoots the current iterate.
\item \textbf{Hyperparameter sensitivity.} The only tunable scalar is $\eta$; any $\eta\le 1/L$ yields the $O(1/k^{2})$ bound. Larger $\eta$ may violate (A3) and break the proof.
\item \textbf{Comparison.}
    \begin{itemize}
        \item Gradient descent with proximal step (no momentum) attains $O(1/k)$.
        \item PAGM improves the rate by a factor $k$, matching the optimal rate for the class of composite convex problems.
    \end{itemize}
\item \textbf{Special cases.}
    \begin{enumerate}
        \item If $g\equiv0$, PAGM reduces to Nesterov’s accelerated gradient method.
        \item If $\beta_k\equiv0$, the algorithm becomes the standard proximal gradient method.
    \end{enumerate}
\end{itemize}

\section*{Preliminaries (Definitions and Lemmas)}
\begin{lemma}[Descent Lemma]\label{lem:descent}
If $\nabla f$ is $L$‑Lipschitz, then for any $u,v\in\R^{d}$
\[
f(v)\le f(u)+\langle\nabla f(u),v-u\rangle+\frac{L}{2}\|v-u\|^{2}. \tag{L1}
\]
\end{lemma}
\begin{lemma}[Three‑point inequality for proximal operator]\label{lem:prox}
For any $z\in\R^{d}$ and any $u\in\R^{d}$,
\[
g(\prox_{\eta g}(z))+\frac{1}{2\eta}\|\prox_{\eta g}(z)-z\|^{2}
\le g(u)+\frac{1}{2\eta}\|u-z\|^{2}
-\frac{1}{2\eta}\|u-\prox_{\eta g}(z)\|^{2}. \tag{L2}
\]
\end{lemma}

\section*{Main Proof (Step‑by‑Step Derivation)}
We prove Theorem~\ref{thm:rate} by constructing a potential function and showing it decreases geometrically.  
All non‑trivial steps are numbered for easy reference.

\noindent\textbf{Notation.} Let $x^\star\in\arg\min F$ and define
\[
A_k:=t_k^{2}\bigl(F(x_k)-F(x^\star)\bigr)+\frac{1}{2\eta}\|x^\star-x_k\|^{2}. \tag{P1}
\]

\noindent\textbf{Goal.} Show $A_{k+1}\le A_k$ for every $k\ge0$.  
Since $t_k\ge k/2$ (standard property of the recurrence), the bound in Theorem~\ref{thm:rate} follows.

\bigskip
\noindent\textbf{Step 1.} Expand $A_{k+1}$ using the definition of $x_{k+1}$.

\[
\begin{aligned}
A_{k+1}
&=t_{k+1}^{2}\bigl(F(x_{k+1})-F(x^\star)\bigr)+\frac{1}{2\eta}\|x^\star-x_{k+1}\|^{2}\\
&=t_{k+1}^{2}\bigl(f(x_{k+1})+g(x_{k+1})-F(x^\star)\bigr)
   +\frac{1}{2\eta}\|x^\star-x_{k+1}\|^{2}. \tag{S1}
\end{aligned}
\]

\noindent\textbf{Step 2.} Apply Lemma~\ref{lem:descent} with $u=y_k$ and $v=x_{k+1}$.

\[
f(x_{k+1})\le f(y_k)+\langle\nabla f(y_k),x_{k+1}-y_k\rangle+\frac{L}{2}\|x_{k+1}-y_k\|^{2}. \tag{S2}
\]

\noindent\textbf{Step 3.} Use the proximal optimality condition (Lemma~\ref{lem:prox}) with $z=y_k-\eta\nabla f(y_k)$ and $u=x^\star$:

\[
\begin{aligned}
g(x_{k+1}) &+\frac{1}{2\eta}\|x_{k+1}-\bigl(y_k-\eta\nabla f(y_k)\bigr)\|^{2}\\
&\le g(x^\star)+\frac{1}{2\eta}\|x^\star-\bigl(y_k-\eta\nabla f(y_k)\bigr)\|^{2}
-\frac{1}{2\eta}\|x^\star-x_{k+1}\|^{2}. \tag{S3}
\end{aligned}
\]

\noindent\textbf{Step 4.} Add inequalities (S2) and (S3).  The terms $\langle\nabla f(y_k),x_{k+1}-y_k\rangle$ cancel because  

\[
\|x_{k+1}-(y_k-\eta\nabla f(y_k))\|^{2}
=\|x_{k+1}-y_k\|^{2}-2\eta\langle\nabla f(y_k),x_{k+1}-y_k\rangle+\eta^{2}\|\nabla f(y_k)\|^{2}.
\]

Substituting this identity into (S3) and adding to (S2) yields

\[
\begin{aligned}
F(x_{k+1})-F(x^\star)
&\le f(y_k)-f(x^\star)+\frac{L}{2}\|x_{k+1}-y_k\|^{2}\\
&\quad+\frac{1}{2\eta}\|x^\star-y_k+\eta\nabla f(y_k)\|^{2}
      -\frac{1}{2\eta}\|x^\star-x_{k+1}\|^{2}\\
&\quad-\frac{1}{2\eta}\|x_{k+1}-y_k\|^{2}
      +\frac{\eta}{2}\|\nabla f(y_k)\|^{2}. \tag{S4}
\end{aligned}
\]

\noindent\textbf{Step 5.} Simplify (S4).  Using $\eta\le 1/L$ we have $\frac{L}{2}-\frac{1}{2\eta}\le0$, thus the term $\|x_{k+1}-y_k\|^{2}$ disappears:

\[
F(x_{k+1})-F(x^\star)
\le f(y_k)-f(x^\star)
   +\frac{1}{2\eta}\bigl\|x^\star-y_k+\eta\nabla f(y_k)\bigr\|^{2}
   -\frac{1}{2\eta}\|x^\star-x_{k+1}\|^{2}. \tag{S5}
\]

\noindent\textbf{Step 6.} Expand the square in (S5):

\[
\begin{aligned}
\bigl\|x^\star-y_k+\eta\nabla f(y_k)\bigr\|^{2}
&=\|x^\star-y_k\|^{2}+2\eta\langle\nabla f(y_k),x^\star-y_k\rangle+\eta^{2}\|\nabla f(y_k)\|^{2}.
\end{aligned}
\]

Insert this into (S5) and rearrange:

\[
\begin{aligned}
F(x_{k+1})-F(x^\star)
&\le f(y_k)-f(x^\star)+\frac{1}{2\eta}\|x^\star-y_k\|^{2}
   -\frac{1}{2\eta}\|x^\star-x_{k+1}\|^{2}\\
&\quad+\langle\nabla f(y_k),x^\star-y_k\rangle+\frac{\eta}{2}\|\nabla f(y_k)\|^{2}. \tag{S6}
\end{aligned}
\]

\noindent\textbf{Step 7.} Use convexity of $f$:
\[
f(x^\star)\ge f(y_k)+\langle\nabla f(y_k),x^\star-y_k\rangle. \tag{S7}
\]
Hence $f(y_k)-f(x^\star)+\langle\nabla f(y_k),x^\star-y_k\rangle\le0$ and (S6) simplifies to

\[
F(x_{k+1})-F(x^\star)
\le \frac{1}{2\eta}\|x^\star-y_k\|^{2}
   -\frac{1}{2\eta}\|x^\star-x_{k+1}\|^{2}
   +\frac{\eta}{2}\|\nabla f(y_k)\|^{2}. \tag{S8}
\]

\noindent\textbf{Step 8.} Multiply (S8) by $t_{k+1}^{2}$ and add $\frac{1}{2\eta}\|x^\star-x_{k+1}\|^{2}$ on both sides:

\[
\begin{aligned}
t_{k+1}^{2}\bigl(F(x_{k+1})-F(x^\star)\bigr)
&+\frac{1}{2\eta}\|x^\star-x_{k+1}\|^{2} \\
&\le t_{k+1}^{2}\Bigl[\frac{1}{2\eta}\|x^\star-y_k\|^{2}
   +\frac{\eta}{2}\|\nabla f(y_k)\|^{2}\Bigr]. \tag{S9}
\end{aligned}
\]

\noindent\textbf{Step 9.} Relate $y_k$ to $x_k$ and $x_{k-1}$.  By definition,
\[
y_k = x_k + \beta_k (x_k-x_{k-1}),\qquad
\beta_k = \frac{t_{k-1}-1}{t_k}. \tag{S10}
\]
A straightforward algebraic computation (see Lemma~\ref{lem:aux}) gives

\[
t_{k+1}^{2}\|x^\star-y_k\|^{2}
\le t_k^{2}\|x^\star-x_k\|^{2}
      - t_{k-1}^{2}\|x^\star-x_{k-1}\|^{2}. \tag{S11}
\]

\begin{lemma}\label{lem:aux}
For the sequence $\{t_k\}$ defined in the algorithm, identity (S11) holds.
\end{lemma}
\begin{proof}
Using $t_{k+1}^{2}=t_k^{2}+t_{k+1}$ (a consequence of the quadratic equation defining $t_{k+1}$) and the definition of $\beta_k$, one expands $t_{k+1}^{2}\|x^\star-y_k\|^{2}$ and simplifies; the terms cancel exactly, yielding (S11).  ∎
\end{proof}

\noindent\textbf{Step 10.} Plug (S11) into (S9).  The term $\frac{\eta}{2}t_{k+1}^{2}\|\nabla f(y_k)\|^{2}$ is non‑positive because $\eta\le 1/L$ and by the descent lemma we have $\|\nabla f(y_k)\|^{2}\le 2L\bigl(f(y_k)-f(x^\star)\bigr)$.  Dropping this non‑negative contribution gives

\[
\begin{aligned}
A_{k+1}
&=t_{k+1}^{2}\bigl(F(x_{k+1})-F(x^\star)\bigr)
   +\frac{1}{2\eta}\|x^\star-x_{k+1}\|^{2}\\
&\le t_k^{2}\bigl(F(x_k)-F(x^\star)\bigr)
   +\frac{1}{2\eta}\|x^\star-x_k\|^{2}
   -\Bigl[t_{k-1}^{2}\bigl(F(x_{k-1})-F(x^\star)\bigr)
   +\frac{1}{2\eta}\|x^\star-x_{k-1}\|^{2}\Bigr] \\
&= A_k - A_{k-1}. \tag{S12}
\end{aligned}
\]

Since $A_{-1}=0$ (by construction) we obtain by induction $A_{k+1}\le A_k$ for all $k\ge0$.

\noindent\textbf{Step 11.} Bounding $t_k$.  From the recursion $t_{k+1}=\frac{1+\sqrt{1+4t_k^{2}}}{2}$ one proves by induction $t_k\ge \frac{k+1}{2}$.  Consequently $t_k^{2}\ge \frac{(k+1)^{2}}{4}$.

\noindent\textbf{Step 12.} Combine $A_k\le A_0$ with the lower bound on $t_k$:

\[
\frac{(k+1)^{2}}{4}\bigl(F(x_k)-F(x^\star)\bigr)
\le A_k\le A_0=\frac{1}{2\eta}\|x_0-x^\star\|^{2}.
\]

Rearranging yields the claimed rate:

\[
F(x_k)-F(x^\star)\le\frac{2\|x_0-x^\star\|^{2}}{\eta\,(k+1)^{2}},\qquad k\ge0. \tag{S13}
\]

This completes the proof of Theorem~\ref{thm:rate}. \hfill $\blacksquare$

\section*{Rate Derivation (With Constants)}
The constant in (S13) is $C:=2\|x_0-x^\star\|^{2}/\eta$.  Therefore the algorithm achieves
\[
\boxed{F(x_k)-F(x^\star)\le \dfrac{C}{k^{2}}},\qquad\text{with }C= \frac{2}{\eta}\|x_0-x^\star\|^{2}.
\]
If $L$ is known, a common choice $\eta=1/L$ gives $C=2L\|x_0-x^\star\|^{2}$.

\section*{Special Cases}
\begin{enumerate}
\item \textbf{Pure smooth minimization ($g\equiv0$).} The proximal step reduces to the identity and the method becomes Nesterov’s accelerated gradient, preserving the $O(1/k^{2})$ rate.
\item \textbf{No momentum ($\beta_k=0$).} Setting $t_k=1$ for all $k$ collapses the algorithm to the proximal gradient method, for which the same proof yields the classical $O(1/k)$ bound.
\item \textbf{Strongly convex $F$.} If $F$ is $\mu$‑strongly convex, a slight modification of the potential (adding $\mu$‑terms) leads to a linear convergence rate $O\bigl((1-\sqrt{\mu/L})^{k}\bigr)$.
\end{enumerate}

\end{document}